\chapter{正则系综}\label{ch:b3:canonical-ensemble}

孤立系统只是理论家的虚构：真实样品置身于
恒温器、房间、空气的海洋之中——与热库接触，
热库固定的不是它们的能量，而是它们的\emph{温度}。本科二年级
分册以经验方式得到了这一规律：能量为
$E$ 的状态的概率带有因子 $\eu^{-E/k_{\text{B}}T}$。本章
本章只用四行，由
\cref{ch:b3:microcanonical-ensemble}的计数推出该玻尔兹曼因子，然后构造
使统计物理成为一门工业学科的机器：
\emph{\omterm{thm:b3:canonical-ensemble:boltzmann}{配分函数}} $Z$，一个求和，由之可得能量、熵、
压强、热容与涨落，皆可
通过求导得到。此处记述这台机器的最初成就：
为何低温下热容会消失（爱因斯坦，1907 年——
物质的第一个量子理论），为何氢的热容
沿台阶攀升，为何天空按指数规律变稀薄——以及
佩兰观察水滴中微小颗粒的沉降，
测得阿伏伽德罗常数，终结了关于
原子是否真实存在的争论。

\section{由此推出的玻尔兹曼分布}

\begin{theorem}[正则分布]\label{thm:b3:canonical-ensemble:boltzmann}
\index{玻尔兹曼分布}\index{正则系综}\index{配分函数}
与温度为 $T$ 的大热库交换能量的系统，
占据其\omterm{def:b3:microcanonical-ensemble:states}{微观态} $s$（能量为 $E_s$）的概率为
\[
\mathcal P_s = \frac{\eu^{-E_s/k_{\text{B}}T}}{Z} , \qquad
Z = \sum_s \eu^{-E_s/k_{\text{B}}T}
\]
（以下记 $\beta = 1/k_{\text{B}}T$。）$Z$，即
\emph{\omterm{thm:b3:canonical-ensemble:boltzmann}{配分函数}}，它使概率归一化，
并恰好包含该系统的全部热力学信息。
\end{theorem}

\begin{proof}
处于态 $s$ 的系统使热库失去能量 $E - E_s$：
由作用于整体的基本假设，$\mathcal P_s
\propto \Omega_{\text{库}}(E - E_s)$。将热库的
熵（巨大而平缓）展开到一阶：$S_{\text{库}}(E - E_s)
= S_{\text{库}}(E) - E_s\,\partial S_{\text{库}}/\partial E =
S_{\text{库}}(E) - E_s/T$，由温度的定义。
因此 $\Omega_{\text{库}} \propto \eu^{-E_s/k_{\text{B}}T}$。（更高阶
的项随热库的尺度而消失。）
\end{proof}

\begin{proposition}[机器]\label{prop:b3:canonical-ensemble:machine}
\index{自由能}\index{热容!与涨落}
由 $Z(T, V, N)$ 可得：
\[
\langle E\rangle = -\frac{\partial\ln Z}{\partial\beta} , \qquad
F = -k_{\text{B}}T\ln Z , \qquad
S = -\frac{\partial F}{\partial T} , \qquad
P = -\frac{\partial F}{\partial V} ,
\]
其中 $F = \langle E\rangle - TS$ 称为\emph{\omterm{prop:b3:canonical-ensemble:machine}{自由能}}；而
热容则是一种\emph{涨落}：
\[
C = \frac{\partial\langle E\rangle}{\partial T}
 = \frac{\langle E^2\rangle - \langle E\rangle^2}
{k_{\text{B}}T^2} .
\]
对独立、可分辨的子系统，$Z = z^N$；对 $N$
个\omterm{def:b3:identical-particles:exchange}{全同粒子}处于同一容器中，$Z = z^N/N!$。在固定的 $T$ 与
$V$ 下，平衡态使 $F$ \emph{取极小}：自然界在能量与
熵之间以温度 $T$ 为兑换率进行权衡——这是
物质物理中最有用的一条原理。
\end{proposition}

\begin{proof}[部分证明]
$-\partial_\beta\ln Z = \sum E_s\eu^{-\beta E_s}/Z = \langle
E\rangle$；再求一次导数得 $\langle E^2\rangle - \langle
E\rangle^2$，由链式法则把 $\partial_\beta$ 换成
$\partial_T$。$F$、$S$、$P$ 的定义由
比较 $\dd(-k_{\text{B}}T\ln Z)$ 与热力学的
关系 $\dd F = -S\dd T - P\dd V$ 得出；因子分解即
一个和的指数。极小原理：$F_{\text{系统}}$ 减小
等价于 $S_{\text{总}}$ 增大，因为 $\Delta S_{\text{库}} =
-\Delta E_{\text{系统}}/T$.
\end{proof}

\begin{omfigure}
\begin{tikzpicture}[font=\small, >=Stealth]
  \fill[omExo!12, draw=omExo, thick, rounded corners=3pt] (-4.4,-1.9) rectangle (4.4,1.9);
  \node[omExo, anchor=north west, font=\footnotesize] at (-4.3,1.82)
    {温度为 $T$ 的热库：固定兑换率};
  \fill[omDef!22, draw=omDef, very thick, rounded corners=2pt] (-1.05,-0.8) rectangle (1.05,0.8);
  \node[omDef] at (0,0.22) {系统};
  \node[omDef, font=\footnotesize] at (0,-0.3) {状态 $s$，$E_s$};
  \draw[<->, thick] (1.15,0) -- (2.5,0);
  \node[anchor=south, font=\footnotesize] at (1.85,0.08) {能量};
  \node[anchor=west, align=left, font=\footnotesize] at (5.0,0.6)
    {$\mathcal P_s \propto \eu^{-E_s/k_{\text{B}}T}$};
  \node[anchor=west, align=left, font=\footnotesize] at (5.0,-0.4)
    {借取 $E_s$ 的代价：\\热库熵 $-E_s/T$};
\end{tikzpicture}

{\small 正则情形：小系统从
大热库借取能量。每借取一焦耳，热库损失 $1/T$ 的
熵——因此高能态受到指数式的折扣。}
\end{omfigure}

\section{最初的成就}

\begin{example}[两能级，两行算完]\label{ex:b3:canonical-ensemble:twolevel}
对能级 $0, \epsilon$，有 $z = 1 + \eu^{-\beta\epsilon}$，于是
$\langle E\rangle = N\epsilon/(\eu^{\beta\epsilon} + 1)$——这一
结果在微正则途径中需要一页斯特林公式的计算
（见\cref{exo:b3:microcanonical-ensemble:8}）。正则形式就是
即预先整理好组合计数的微正则方法；
$C(T)$ 中的肖特基峰只需再求一次导数即可得到。
\end{example}

\begin{example}[爱因斯坦固体与杜隆–珀蒂定律的失效]\label{ex:b3:canonical-ensemble:einstein}
\index{爱因斯坦固体}\index{杜隆–珀蒂定律}
将晶体视为 $3N$ 个频率为 $\omega$ 的量子振子，
（见\cref{ch:b3:harmonic-oscillator}）。对每个振子，
\[
z = \sum_n\eu^{-\beta\hbar\omega(n + 1/2)}
 = \frac{1}{2\sinh(\beta\hbar\omega/2)} , \qquad
\langle E\rangle = \frac{\hbar\omega}{2} +
\frac{\hbar\omega}{\eu^{\beta\hbar\omega} - 1} .
\]
在高温极限 $T$ 下，每个振子 $\langle E\rangle \to k_{\text{B}}T$，
$C \to 3Nk_{\text{B}}$——本科一年级分册中的杜隆–珀蒂定律，得到解释。
低 $T$ 时，量子 $\hbar\omega$ 变得
难以被热能激发，$C$ 指数式骤降——与实测相符，而
经典理论无法解释。爱因斯坦 1907 年的曲线，每种元素仅一个参数，
是量子概念首次用于普通物质；
金刚石键刚度大、原子质量轻（$\hbar\omega/k_{\text{B}}
\approx \qty{1300}{K}$），室温下仍处于“冻结”状态——
这一曾困扰化学家八十年的反常
（见\cref{exo:b3:canonical-ensemble:6}）。
\end{example}

\begin{omfigure}
\begin{tikzpicture}
  \begin{axis}[omaxis, axis x line=bottom, axis y line=left,
    width=6.8cm, height=5.2cm, xmin=0, xmax=2, ymin=0, ymax=1.15,
    xlabel={$T/\theta_{\text{E}}$}, ylabel={$C/3Nk_{\text{B}}$},
    xlabel style={at={(axis description cs:0.5,-0.11)}, anchor=north},
    xtick={1,2}, ytick={1}, clip=false]
    \addplot[gray, dashed, domain=0:2] {1};
    \addplot[omDef, very thick, domain=0.02:2, samples=200]
      {(1/x)^2*exp(1/x)/(exp(1/x)-1)^2};
    \node[font=\scriptsize, anchor=west, gray] at (axis cs:0.08,1.07)
      {杜隆–珀蒂};
    \node[omDef, font=\scriptsize, anchor=west] at (axis cs:0.75,0.5)
      {爱因斯坦 1907};
  \end{axis}
  \begin{axis}[omaxis, axis x line=bottom, axis y line=left,
    width=6.8cm, height=5.2cm, xmin=1, xmax=4000, ymin=0, ymax=4.3,
    xshift=7.4cm, xmode=log,
    xlabel={$T$（K，对数刻度）}, ylabel={$C_V/Nk_{\text{B}}$（H$_2$ 气体）},
    xlabel style={at={(axis description cs:0.5,-0.11)}, anchor=north},
    xtick={10,100,1000}, xticklabels={10,100,1000},
    ytick={1.5,2.5,3.5}, clip=false]
    \addplot[omThm, very thick, domain=2:4000, samples=400]
      {1.5 + (87/x)^2*exp(-87/x)/(1-exp(-87/x))^2
       + ((6000/x)<25 ? (6000/x)^2*exp(-6000/x)/(1-exp(-6000/x))^2 : 0)};
    \node[omThm, font=\scriptsize, anchor=west] at (axis cs:3,1.68) {平动};
    \node[omThm, font=\scriptsize, anchor=west] at (axis cs:150,2.72) {+ 转动};
    \node[omThm, font=\scriptsize, anchor=south east] at (axis cs:3800,3.55) {+ 振动};
  \end{axis}
\end{tikzpicture}

{\small 左：爱因斯坦热容曲线——经典能量均分
在高 $T$ 下重新成立，低于
$\theta_{\text{E}} = \hbar\omega/k_{\text{B}}$ 时量子冻结。右：氢气的
$C_V$ 阶梯式上升：先是转动（约 \qty{85}{K}）
随后是振动（约 \qty{6000}{K}，超出量程）解冻：每种运动
只有当 $k_{\text{B}}T$ 能够支付其量子时，才加入能量均分。}
\end{omfigure}

\begin{theorem}[能量均分定理及其适用条件]\label{thm:b3:canonical-ensemble:equipartition}
\index{能量均分定理}
凡以\emph{二次形式}进入能量的每个坐标或动量，
在经典（高温）区域中，其对
\[
\langle\epsilon\rangle = \tfrac12 k_{\text{B}}T
\]
平均能量的贡献为：单原子气体的原子 $\tfrac32 k_{\text{B}}T$，
转动双原子分子为 $\tfrac52$，振子为 $3k_{\text{B}}T$。
适用条件在 $k_{\text{B}}T$ 低于该模式的
能级间距时失效：该模式冻结，其贡献
为零——这便解决了十九世纪的
热容疑难，如上图阶梯所示。
\end{theorem}

\begin{proof}[部分证明]
设 $\epsilon = ax^2$，
\[
\langle\epsilon\rangle
 = \frac{\int ax^2\,\eu^{-\beta ax^2}\dd x}
        {\int\eu^{-\beta ax^2}\dd x}
 = -\partial_\beta\ln\!\int\eu^{-\beta ax^2}\dd x
 = -\partial_\beta\ln\beta^{-1/2} = \frac{1}{2\beta} .
\] 冻结过程即
\cref{ex:b3:canonical-ensemble:einstein}中的计算，逐个
模式分别进行。
\end{proof}

\begin{example}[正则系综下的气体]\label{ex:b3:canonical-ensemble:gas}
箱中一个原子：$z = V/\lambda_T^3$（即
\cref{prop:b3:schrodinger-three-dimensions:counting}中的态数，以玻尔兹曼因子加权）；
$N$ 个全同原子：$Z = z^N/N!$。于是 $F =
-Nk_{\text{B}}T[\ln(V/N\lambda_T^3) + 1]$，求导即得
$PV = Nk_{\text{B}}T$，萨克尔–泰特罗德熵，以及
$\langle E\rangle = \tfrac32 Nk_{\text{B}}T$——整个理想气体
仅由一次高斯积分即可得出。将玻尔兹曼因子用于分子的
动能，即得本科一年级分册中的麦克斯韦速率分布（现已推导）；
用于其势能
$mgh$ 得指数型大气——而在一滴水中，
可见佩兰的颗粒阶梯
（见\cref{pb:b3:canonical-ensemble:1}）。
\end{example}

\begin{method}[正则计算技巧]\label{met:b3:canonical-ensemble:method}
(1) 列出\emph{一个}单元的态与能量；计算 $z$。
(2) 因式分解：$Z = z^N$（对共享空间的\omterm{def:b3:identical-particles:exchange}{全同粒子}则为
$z^N/N!$）；尽早取 $\ln$。(3) 求导：$\beta$ 给出能量，
$T$ 经由 $F$ 求熵，$V$ 求压强；二阶导数
用于 $C$ 与涨落。(4) 检验两端：高 $T$
时必须再现能量均分；低 $T$ 时须冻结，呈
$\eu^{-\Delta/k_{\text{B}}T}$ 形尾部。(5) 竞争（折叠、
结合、取向）：对每种情形写出 $F = E - TS$，并
取较小者胜出——交叉点位于 $T \approx \Delta
E/\Delta S$ 处。
\end{method}

\section{练习}

\begin{exercise}[$\star$]\label{exo:b3:canonical-ensemble:1}
(a) 写出被 $\epsilon$ 分开的两能级
在温度 $T$ 下的粒子数之比。(b) 在 \qty{2500}{K} 火焰中，
钠原子中处于 D 线的 \qty{2.1}{eV} 激发态的比例是多少，
（基态简并度为 2，激发态简并度为 6，粒子数之比为
$3\eu^{-\beta\epsilon}$）？
(c) 火焰为何仍然呈黄色（每
cm$^3$ 需多少原子才足够）？(d) 温度为多少时，激发态占比会达到
$10\%$？
\end{exercise}

\begin{exercise}[$\star$]\label{exo:b3:canonical-ensemble:2}
一个三能级体系：能级为 $0, \epsilon, 2\epsilon$。(a) 写出 $z$。(b)
计算 $\langle E\rangle$，并检验两个温度极限。(c)
在何 $T$ 时中间能级的\emph{绝对
粒子数}最多？(d) 证明：无论温度多高，
较高能级永远不会比较低能级更多地被占据——这需要用到第 16 章的哪个
概念？
\end{exercise}

\begin{exercise}[$\star$]\label{exo:b3:canonical-ensemble:3}
等温大气。(a) 将玻尔兹曼因子用于
势能 $mgh$，并导出 $n(h) = n_0\eu^{-mgh/k_{\text{B}}
T}$。(b) 计算空气（$m = \qty{4.8e-26}{kg}$）的标高
在 \qty{288}{K} 下。(c) 珠穆朗玛峰峰顶压强占
海平面压强的比例。(d) 为何真实大气的
衰减接近指数规律，却并不严格（我们假定为恒定的量中，哪一个实际上并不恒定？）
\end{exercise}

\begin{exercise}[$\star$]\label{exo:b3:canonical-ensemble:4}
能量均分的核算。数清二次项，并预测
以下物质的摩尔热容 $C_V$：(a) 氩；(b) 室温下的 N$_2$
（转动开启、振动冻结）；(c) N$_2$ 在 \qty{3000}{K} 时；(d) 一种
晶体固体（杜隆–珀蒂）。实测值
$12.5$、$20.8$、$\approx26$、$\approx\qty{25}{J/(mol.K)}$ 中哪些吻合，
各偏差又说明了什么？
\end{exercise}

\begin{exercise}[$\star\star$]\label{exo:b3:canonical-ensemble:5}
\omterm{thm:b3:canonical-ensemble:boltzmann}{正则系综}下的量子振子。(a) 对几何级数求和，
得 $z$。(b) 导出 $\langle E\rangle$，并辨认 $\langle
n\rangle = 1/(\eu^{\beta\hbar\omega} - 1)$——正是
\cref{exo:b3:harmonic-oscillator:6}中的结果，如今轻而易举。(c)
对 $C(T)$ 求导，并检验两个极限。(d) 为何
$\langle n\rangle$ 的形式即将变得著名
（见\cref{ch:b3:photons-phonons}）？
\end{exercise}

\begin{exercise}[$\star\star$]\label{exo:b3:canonical-ensemble:6}
爱因斯坦与数据的对比。(a) 由图中公式，在何
$T/\theta_{\text{E}}$ 处 $C$ 降为杜隆–珀蒂值的一半？(b)
金刚石：$\theta_{\text{E}} \approx \qty{1300}{K}$——计算
在 \qty{300}{K} 时的 $C/3Nk_{\text{B}}$，并解释十九世纪的
“金刚石反常”。(c) 铅：$\theta_{\text{E}}
\approx \qty{90}{K}$——为何铅一直“表现良好”？(d)
在极低 $T$ 下的测量表明 $C \propto T^3$，而非
指数形式：爱因斯坦哪一条假设不成立（所有
振子都具有\emph{单一}频率），并由谁修正了这一假设
（见\cref{ch:b3:photons-phonons}）？
\end{exercise}

\begin{exercise}[$\star\star$]\label{exo:b3:canonical-ensemble:7}
麦克斯韦的尾部与缺失的氢。(a) 由作用于动能的玻尔兹曼
因子出发，写出速率分布，并定出
N$_2$ 与 H$_2$ 在 \qty{288}{K} 时的最概然速率。(b)
地球逃逸速率为 \qty{11.2}{km/s}：计算
两种气体的 $mv_{\text{逃逸}}^2/2k_{\text{B}}T$。(c) 逃逸
比例正比于 $\eu^{-mv_{\text{逃逸}}^2/2k_{\text{B}}T}$：比较
两个指数，并断定哪种气体在地质
时间尺度上会逸散。(d) 联系地球大气的实测成分（不含游离
H$_2$），与木星（主要为 H$_2$）大气相比较——哪两个
参数会颠倒这一结论？
\end{exercise}

\begin{exercise}[$\star\star$]\label{exo:b3:canonical-ensemble:8}
涨落与响应。(a) 证明 $C = (\langle E^2\rangle -
\langle E\rangle^2)/k_{\text{B}}T^2$，由 $\ln
Z$ 的两次求导得出。(b) 在两能级体系上明确验证。(c) 对 $N$ 个
独立单元，证明\emph{相对}能量涨落
以 $1/\sqrt N$ 的规律减小。(d) 点明其中的道理：
系统愿意吸收热量（一种响应）与其能量
抖动的幅度（一种涨落）之间有何关系——统计物理
中反复出现的一笔权衡。
\end{exercise}

\begin{exercise}[$\star\star$]\label{exo:b3:canonical-ensemble:9}
化学实验室中的玻尔兹曼。反应速率带有因子
$\eu^{-E_a/k_{\text{B}}T}$（跨越活化能垒 $E_a$
——阿伦尼乌斯）。(a) 证明经验规则“室温附近
每升高 \qty{10}{K} 速率加倍”对应于 $E_a \approx
\qty{0.55}{eV}$。(b) 该反应在冰箱（\qty{5}{\celsius}）中
放慢了多少倍？(c) 高海拔煮蛋：
在 \qty{2000}{m} 高的山口，水在 \qty{93}{\celsius} 沸腾——
估算额外增加的烹饪时间。(d) 为何决定化学
（哪些分子发生反应）的是势垒尾部，而非平均能量？
\end{exercise}

\begin{exercise}[$\star\star\star$]\label{exo:b3:canonical-ensemble:10}
两态折叠。生物分子可处于折叠态（能量 $0$，仅有
一种构型），也可处于展开态（能量 $\Delta E > 0$，$\Omega_u =
\eu^{\Delta S/k_{\text{B}}}$ 种构型）。(a) 写出折叠态的
占比与 $T$ 的关系。(b) 证明“熔化”中点温度为
$T_{\text{m}} = \Delta E/\Delta S$，并将其解释为 $F = E -
TS$ 的相持点。(c) 取 $\Delta E = \qty{3.0}{eV}$ 与 $\Delta S =
100\,k_{\text{B}}$（一个小蛋白质的协同单元）：计算
$T_{\text{m}}$ 及转变的宽度。(d) 为何
\emph{协同性}（许多接触同时断开，
$\Delta E$ \emph{和} $\Delta S$ 均较大）使熔化变得陡峭——而
DNA 热循环仪（PCR）如何恰恰利用这一点？
\end{exercise}

\begin{exercise}[$\star\star\star$]\label{exo:b3:canonical-ensemble:11}
顺磁性与磁制冷。$N$ 个自旋 $\tfrac12$、磁矩为
$\mu$ 的粒子处于磁场 $B$ 中。(a) 证明 $M = N\mu\tanh(\mu B/k_{\text{B}}T)$
并展开得到居里定律 $M \approx N\mu^2B/k_{\text{B}}T$。(b)
对电子磁矩求取向参数 $\mu B/k_{\text{B}}T$ 的值，
分别在 $B = \qty{1}{T}$、$T = \qty{300}{K}$ 与 \qty{1}{K} 下计算。(c)
绝热退磁：在 \qty{1}{K} 下磁化，隔离系统，
将 $B$ 减小为原来的十分之一——证明熵不变意味着
$\mu B/k_{\text{B}}T$ 保持不变，从而 $T$ 降为原来的十分之一。(d) 这一制冷机
的极限温度由什么决定（自旋之间的相互作用——
估算偶极相互作用标度 $\mu_0\mu^2/4\pi a^3$，取 $a =
\qty{0.5}{nm}$，并以温度为单位表示）？
\end{exercise}

\begin{exercise}[$\star\star\star$]\label{exo:b3:canonical-ensemble:12}
转动阶梯的定量分析。双原子分子的转动
能级为 $E_J = BJ(J+1)$，简并度为 $2J + 1$
（见\cref{ch:b3:quantum-angular-momentum}）。(a) 写出 $z_{\text{转动}}$，
并证明当 $k_{\text{B}}T \gg B$ 时，求和化为
积分 $k_{\text{B}}T/B$：能量均分定理的 $k_{\text{B}}$ 在
$C_V$ 中得以恢复。(b) 定义 $\theta_{\text{转动}} =
B/k_{\text{B}}$，并对 H$_2$（$B = \qty{7.5}{meV}$）与
N$_2$（$B = \qty{0.25}{meV}$）求值：哪种气体
在可达温度下出现转动台阶？(c) 画出氢气完整的 $C_V(T)$
阶梯图，标出其三个平台与两段跃升，并
在两者上标注数值。(d) 低 $T$ 时 H$_2$ 的实测行为，还
因 3:1 的正仲氢混合物而更为复杂（见
\cref{exo:b3:identical-particles:8}）：用一句话说明核自旋统计如何
影响气体的热容。
\end{exercise}

\section{问题：在一滴水中数出阿伏伽德罗常数}

\begin{problem}[{周末问题——佩兰的颗粒与
原子的实在性}]\label{pb:b3:canonical-ensemble:1}
1908 年，让·佩兰将微小的树脂颗粒悬浮于水中，
使其沉降，并在显微镜下逐层计数。他发现的指数
阶梯正是玻尔兹曼因子的可见形式——由其标高，
他求得了阿伏伽德罗常数，
使最后的怀疑者相信原子确实存在。
诺贝尔奖，1926 年。数据：藤黄颗粒半径 $r =
\qty{0.212}{\micro m}$，密度 $\rho_{\text{g}} =
\qty{1207}{kg/m^3}$；水的密度 $\rho_{\text{w}} = \qty{999}{kg/m^3}$；
$T = \qty{293}{K}$；$g = \qty{9.81}{m/s^2}$。

\medskip\noindent\textbf{第一部分——可见的玻尔兹曼因子。}
\begin{enumerate}
  \item 水中的颗粒受重力与浮力共同作用：计算其
        \emph{有效}质量 $m' = \tfrac43\pi r^3(\rho_{\text{g}}
        - \rho_{\text{w}})$ 与重量。
  \item 写出由玻尔兹曼因子给出的
        平衡浓度分布 $n(h)$。
  \item 计算标高 $h_0 = k_{\text{B}}T/m'g$，
        取现代的 $k_{\text{B}}$ 值。
  \item 为何颗粒必须大小高度一致（
        $h_0$ 如何依赖于 $r$）？
  \item 将 $h_0$ 与空气分子的同一公式相比较：
        为何颗粒的“大气层”仅有微米量级，
        而空气的却有千米量级？
  \item 佩兰在一次实验中测得相对浓度
        $100 : 55 : 30 : 17$，位于四个等间距的深度处，
        相隔 \qty{30}{\micro m}：验证其为一个指数阶梯，
        并求出其 $h_0$。
\end{enumerate}

\medskip\noindent\textbf{第二部分——称量看不见的东西。}
\begin{enumerate}[resume]
  \item 反过来：由第 6 题测得的 $h_0$ 与已知的
        $m'g$，求出 $k_{\text{B}}$。
  \item 气体常数 $R = \qty{8.314}{J/(mol.K)}$ 早已由宏观化学得知，
        将其与你求得的 $k_{\text{B}}$ 结合，
        得到阿伏伽德罗常数 $N_{\text{A}} =
        R/k_{\text{B}}$。
  \item 佩兰的实验给出 $N_{\text{A}}$ 介于 $5.5$ 与
        $\num{7.2e23}$ 之间：与现代值
        \num{6.022e23} 比较，并结合他所用的
        显微镜与秒表，评述这一成就。
  \item 由 $N_{\text{A}}$ 计算单个氢原子的
        质量——这是原子论者渴求了一个
        世纪之久的数值。
  \item 解释其逻辑结构：哪些\emph{宏观}
        测量（$R$；颗粒的大小与密度；阶梯分布）
        结合起来，可以在从未见过原子的情况下称出单个原子的质量？
  \item 为何这些颗粒——每个约含 $10^{10}$ 个原子质量——
        服从与分子相同的玻尔兹曼统计（
        在推导中，
        \cref{thm:b3:canonical-ensemble:boltzmann}的哪一步依赖于
        粒子大小）？
\end{enumerate}

\medskip\noindent\textbf{第三部分——确证结论的抖动。}
\begin{enumerate}[resume]
  \item 同一批颗粒还会抖动：爱因斯坦 1905 年关于
        布朗运动的公式给出 $\langle x^2\rangle = 2Dt$，其中 $D =
        k_{\text{B}}T/6\pi\eta r$（水：$\eta =
        \qty{e-3}{Pa.s}$）。计算佩兰颗粒的 $D$。
  \item 一分钟内颗粒会漂移多远？佩兰能否
        用测微目镜与秒表测得它？
  \item 佩兰验证了 $\langle x^2\rangle \propto t$，并
        独立地\emph{再次}求出了 $k_{\text{B}}$：
        为何两条互不相关的途径（静态的阶梯分布；动态的
        抖动）得到同一数值，具有如此重要的证据价值？
  \item 抖动是能量均分定理在颗粒上的应用：每个
        速度分量携带 $\tfrac12 k_{\text{B}}T$。
        估算颗粒的方均根热速率（$m \approx
        \qty{4.8e-17}{kg}$）。
  \item 为何这一速率从未被直接看到（什么打断了
        纳米量级的自由飞行之后），而\emph{实际看到的}
        是什么？
  \item 指出在这一观测中相遇的本书两章：
        阻力的力学（本科二年级分册，黏性）与
        本章的统计学。
\end{enumerate}

\medskip\noindent\textbf{第四部分——已确立的结论。}
\begin{enumerate}[resume]
  \item 奥斯特瓦尔德和马赫曾认为原子只是记账用的虚构：
        用一句话说明，为什么由晶粒阶梯\emph{数出}的
        $N_{\text{A}}$ 终结了这一立场。
  \item 列出 1900 年代另外三条汇聚到同一
        $N_{\text{A}}$ 的途径（天空的蓝色，
        \cref{exo:b3:microcanonical-ensemble:12}；电解加电子电荷；
        放射性产生氦
        ）——为什么\emph{汇聚}
        比任何单一数值更重要？
  \item 佩兰的阶梯是本章哪两种
        “货币”之间的平衡（能量下拉，熵向上扩散），
        以何种比率定价？
  \item 同一物理的现代应用：分析型
        超速离心机以 $10^5g$ 的加速度旋转蛋白质，使其
        “大气”压缩成可测量的阶梯——证明
        把 $g$ 乘以 $10^5$ 会使 $h_0$ 除以同一
        因子，并估算有效质量为
        \qty{e-22}{kg} 的蛋白质的 $h_0$，取 $10^5g$、\qty{293}{K}。
  \item 将晶粒半径加倍会使 $h_0$ 除以八：对于深度为
        \qty{100}{\micro m} 的显微镜视场，给出晶粒尺寸的实用区间，
        其界限分别为“阶梯太高以至于看不到梯度”
        和“阶梯薄于一个晶粒”。
  \item 自 2019 年起，国际单位制\emph{精确定义}了 $k_{\text{B}}$ 和
        $N_{\text{A}}$，说明一次精确的佩兰式实验
        \emph{如今}测量的是什么（一致性检验，或
        对仪器和晶粒的标定），并说明为什么
        物理本身并未改变。
  \item 概括这一具名结果：直径 \qty{0.2}{\micro m} 的晶粒，
        其浓度每升高 $\sim\qty{35}{\micro m}$ 减半；
        用 $n(h) = n_0\eu^{-m'gh/k_{\text{B}}T}$ 读取，
        该阶梯得出 $N_{\text{A}} \approx \num{6e23}$ ——
        原子是被数出的，而非被猜测的，就在
        显微镜载物台上的一滴水中。
\end{enumerate}
\end{problem}
